\section*{Hoofdstuk \ref{ch:b3:point-groups} --- Symmetrie en puntgroepen}

\begin{solution}{exo:b3:point-groups:1}
Etheen: drie loodrechte $C_2$, drie vlakken, $i$: $D_{2h}$. \ce{XeF4} (vlak
vierkant): $C_4$, vier $C_2 \perp C_4$, $\sigma_h$, $2\sigma_v$, $2\sigma_d$, $i$, $S_4$:
$D_{4h}$. \ce{PCl5} (trigonale bipiramide): $C_3$, drie $C_2$, $\sigma_h$, $3\sigma_v$,
$S_3$: $D_{3h}$. 1,3,5-Trichloorbenzeen: dezelfde elementen als \ce{BF3}: $D_{3h}$.
\end{solution}

\begin{solution}{exo:b3:point-groups:2}
Gestaffeld ethaan: $C_3$, drie loodrechte $C_2$, drie $\sigma_d$, $i$, $S_6$:
$D_{3d}$. Eclips: $C_3$, drie $C_2$, $\sigma_h$, drie $\sigma_v$: $D_{3h}$.
Cyclohexaan in de stoelvorm: $D_{3d}$.
\end{solution}

\begin{solution}{exo:b3:point-groups:3}
\ce{SO3} ($D_{3h}$): nee. \ce{SO2} ($C_{2v}$): ja. \ce{PCl3} ($C_{3v}$): ja.
\ce{XeF4} ($D_{4h}$): nee. \ce{CH2Cl2} ($C_{2v}$): ja. cis-1,2-dichlooretheen
($C_{2v}$): ja; trans ($C_{2h}$): nee.
\end{solution}

\begin{solution}{exo:b3:point-groups:4}
$C_2(z) = \mathrm{diag}(-1,-1,1)$, $\chi = -1$; $i = \mathrm{diag}(-1,-1,-1)$,
$\chi = -3$; $\sigma_h(xy) = \mathrm{diag}(1,1,-1)$, $\chi = 1$.
\end{solution}

\begin{solution}{exo:b3:point-groups:5}
Met de diagonaalmatrices van de vorige oefening: $C_2i = \sigma_h$, $C_2\sigma_h =
i$, $i\sigma_h = C_2$, elke operatie is haar eigen inverse, en elk product
is commutatief. De groep is commutatief, dus $X^{-1}AX = A$ voor alle $X$:
elke operatie is een klasse op zich (vier klassen, vier \omterm{def:b3:point-groups:irrep}{irreducibele
representaties} van dimensie 1).
\end{solution}

\begin{solution}{exo:b3:point-groups:6}
Neem voor $\sigma_v(1)$ het $xz$-vlak. Het punt $(1,0)$ gaat door $C_3$ naar $(-\frac12,
\frac{\sqrt3}2)$, en daarna door $\sigma_v(1)$ naar $(-\frac12,-\frac{\sqrt3}2)$; in de
andere volgorde gaat het naar $(1,0)$ en daarna naar $(-\frac12,\frac{\sqrt3}2)$. De twee
producten zijn spiegelingen in twee verschillende vlakken ($\sigma_v(3)$ en
$\sigma_v(2)$): de groep is niet commutatief, en daarom vormen haar
spiegelingen een klasse van drie.
\end{solution}

\begin{solution}{exo:b3:point-groups:7}
Gedraaid bifenyl behoudt de $C_2$ langs de binding tussen de ringen en twee
$C_2$ loodrecht daarop, die de hoeken tussen de ringvlakken middendoor
delen; alle vlakken en het \omterm{def:b3:point-groups:inversion}{inversiecentrum} van de vlakke of loodrechte
vormen gaan verloren. Met drie loodrechte $C_2$ en geen oneigenlijk element
is de groep
$D_2$: het molecuul is chiraal (zijn twee gedraaide vormen zijn
enantiomeren, scheidbaar wanneer omvangrijke orthosubstituenten de rotatie
verhinderen).
\end{solution}

\begin{solution}{exo:b3:point-groups:8}
$d_{z^2}$ en $d_{x^2-y^2}$: $a_1$ ($x^2$, $y^2$, $z^2$ zijn $A_1$); $d_{xy}$: $a_2$; $d_{xz}$:
$b_1$; $d_{yz}$: $b_2$.
\end{solution}

\begin{solution}{exo:b3:point-groups:9}
$1 + 1 + 4 + 9 + 9 = 24 = h$. $\sum g\,\chi_{A_1}\chi_{T_2} = 1\cdot3 + 8\cdot0 + 3(-1)
+ 6(-1) + 6(1) = 0$.
\end{solution}

\begin{solution}{exo:b3:point-groups:10}
In het vlak loodrecht op de snijlijn beeldt een spiegeling in een rechte
onder de hoek $\alpha$ de poolhoek $\phi$ af op $2\alpha - \phi$. Twee
spiegelingen, onder
$\alpha$ en daarna onder $\alpha + \theta$: $\phi \to 2\alpha - \phi \to 2(\alpha + \theta) - (2\alpha - \phi) =
\phi + 2\theta$, een rotatie over $2\theta$ om de snijlijn. Twee verticale
vlakken onder
$60^\circ$ brengen dus een rotatie over $120^\circ$ voort: een $C_3$.
\end{solution}

\begin{solution}{exo:b3:point-groups:11}
De as C=C=C is een $C_2$ en een $S_4$ (draai over $90^\circ$, wat de twee
\ce{CH2}-vlakken verwisselt, en spiegel dan in het vlak door het centrale
koolstofatoom); de twee
\ce{CH2}-vlakken zijn $\sigma_d$; twee $C_2$ loodrecht op de as delen ze
middendoor.
Orde 8: $D_{2d}$, achiraal. In \ce{ClHC=C=CHCl} worden de vlakken, de $S_4$
en de axiale
$C_2$ door de substitutie vernietigd; alleen één $C_2$ loodrecht op de
as blijft over: $C_2$, chiraal --- een axiaal chiraal cumuleen.
\end{solution}

\begin{solution}{exo:b3:point-groups:12}
\ce{SF6}: $O_h$ ($h = 48$). \ce{SF5Cl}: de $C_4$ door Cl en vier $\sigma_v$:
$C_{4v}$ ($h = 8$). trans-\ce{SF4Cl2}: $D_{4h}$ ($h = 16$). cis-\ce{SF4Cl2}: $C_{2v}$ ($h
= 4$). Elk behoudt een deelverzameling van de operaties van $O_h$ die
gesloten is onder producten,
en 8, 16, 4 zijn delers van 48. Polair: \ce{SF5Cl} en cis-\ce{SF4Cl2}.
\end{solution}

\begin{solution}{pb:b3:point-groups:1}
\textbf{1.} Langs de vier lichaamsdiagonalen door C en elk H; elk geeft
$C_3$ en $C_3^2$: 8 rotaties.
\textbf{2.} Langs $x$, $y$, $z$ door de middelpunten van de zijvlakken: $C_2(z)$ beeldt $(1,1,1)
\to (-1,-1,1)$ af, een waterstofatoom. Een rotatie over $90^\circ$ om $z$
gevolgd door een spiegeling in $xy$ beeldt $(1,1,1) \to (-1,1,1) \to (-1,1,-1)$ af, een
waterstofatoom: een
$S_4$; met $S_4^3$ erbij, 6 operaties.
\textbf{3.} De vlakken $x = \pm y$, $y = \pm z$, $x = \pm z$; $x = y$ bevat
$(1,1,1)$ en $(-1,-1,1)$: elk vlak bevat twee bindingen C--H.
\textbf{4.} $1 + 8 + 3 + 6 + 6 = 24$.
\textbf{5.} Inversie beeldt $(1,1,1)$ af op $(-1,-1,-1)$, en dat is geen
waterstofatoom: geen $i$.
\textbf{6.} $E$; $8C_3$; $3C_2$; $6S_4$; $6\sigma_d$.
\textbf{7.} \ce{CH3Cl}: $C_{3v}$, $h = 6$.
\textbf{8.} \ce{CH2Cl2}: $C_{2v}$, $h = 4$.
\textbf{9.} \ce{CHCl3}: $C_{3v}$; \ce{CCl4}: $T_d$, $h = 24$.
\textbf{10.} \ce{CH2FCl}: $C_s$ ($h = 2$); \ce{CHFClBr}: $C_1$ ($h = 1$).
\textbf{11.} Elk behoudt de operaties van $T_d$ die de substitutie
respecteren; 6, 4, 6, 24, 2, 1 zijn alle delers van 24.
\textbf{12.} De drie waterstofatomen, gepermuteerd door $C_3$, $C_3^2$ en de
drie $\sigma_v$.
\textbf{13.} Alle behalve \ce{CCl4}: langs de $C_3$-as voor \ce{CH3Cl} en
\ce{CHCl3}, langs de $C_2$-as voor \ce{CH2Cl2}, in het \omterm{def:b3:point-groups:mirror}{spiegelvlak} voor
\ce{CH2FCl}, in geen opgelegde richting voor \ce{CHFClBr}.
\textbf{14.} Alleen \ce{CHFClBr} ($C_1$, geen $S_n$).
\textbf{15.} Twee enantiomeren.
\textbf{16.} Zijn twee waterstofatomen zijn equivalent: het vlak door C, F
en Cl dat de hoek H--C--H middendoor deelt, is een \omterm{def:b3:point-groups:mirror}{spiegelvlak}.
\textbf{17.} Nee: vier verschillende substituenten laten alleen $E$ over;
zonder $S_n$ is het molecuul chiraal (en kan het polair zijn).
\textbf{18.} $(x, y, z)$: $T_2$.
\textbf{19.} Twee waterstofatomen zijn de uiteinden van een ribbe van de
tetraëder; de 24 operaties brengen elke ribbe op elke andere (de 6 ribben
vormen één verzameling): één \ce{CH2Cl2}.
\textbf{20.} \ce{CH2FCl}: één. \ce{CHFClBr}: twee, een paar enantiomeren.
\textbf{21.} Drie: ortho ($C_{2v}$), meta ($C_{2v}$), para ($D_{2h}$).
\textbf{22.} Drie: 1,2,3- ($C_{2v}$), 1,2,4- ($C_s$), 1,3,5- ($D_{3h}$).
\textbf{23.} Twee (Cl-atomen naast elkaar of tegenover elkaar). Er is maar
één \ce{CH2Cl2} bekend, wat een vlak koolstofatoom uitsloot en het
tetraëdrische koolstofatoom steunde dat in 1874 werd voorgesteld.
\textbf{24.} $1^2 + 1^2 + 2^2 + 3^2 + 3^2 = 24$: \textbf{de orde van $T_d$, de
\omterm{def:b3:point-groups:point-group}{puntgroep} van methaan, is $h = 24$}.
\end{solution}
